\documentclass[UTF8,a4paper,11pt]{ctexart}

\usepackage[a4paper,margin=2.4cm]{geometry}
\usepackage{graphicx}
\usepackage{booktabs}
\usepackage{tabularx}
\usepackage{longtable}
\usepackage{array}
\usepackage{multirow}
\usepackage{xcolor}
\usepackage{enumitem}
\usepackage{hyperref}
\usepackage{amsmath,amssymb}
\usepackage{tikz}
\usepackage[most]{tcolorbox}
\usepackage{titlesec}
\usepackage{fancyhdr}
\usepackage{caption}
\usepackage{float}
\providecommand{\tightlist}{\setlength{\itemsep}{0pt}\setlength{\parskip}{0pt}}
\newcounter{none}

\usetikzlibrary{positioning,arrows.meta,shapes.geometric,fit,calc}

\definecolor{mainblue}{RGB}{28,90,160}
\definecolor{lightblue}{RGB}{235,244,255}
\definecolor{darkgray}{RGB}{70,70,70}
\definecolor{lightgray}{RGB}{245,245,245}
\definecolor{accentorange}{RGB}{230,120,40}
\definecolor{softgreen}{RGB}{235,250,240}

\hypersetup{
  colorlinks=true,
  linkcolor=mainblue,
  urlcolor=mainblue,
  citecolor=mainblue
}

\titleformat{\section}
  {\Large\bfseries\color{mainblue}}
  {\thesection}{0.8em}{}

\titleformat{\subsection}
  {\large\bfseries\color{darkgray}}
  {\thesubsection}{0.8em}{}

\titleformat{\subsubsection}
  {\normalsize\bfseries\color{darkgray}}
  {\thesubsubsection}{0.8em}{}

\pagestyle{fancy}
\setlength{\headheight}{14pt}
\fancyhf{}
\lhead{选题前置调研}
\rhead{研究方案报告}
\cfoot{\thepage}

\tcbset{
  mybox/.style={
    colback=lightblue,
    colframe=mainblue,
    arc=2mm,
    boxrule=0.8pt,
    left=2mm,
    right=2mm,
    top=1.5mm,
    bottom=1.5mm
  },
  keybox/.style={
    colback=softgreen,
    colframe=green!50!black,
    arc=2mm,
    boxrule=0.8pt,
    left=2mm,
    right=2mm,
    top=1.5mm,
    bottom=1.5mm
  },
  warnbox/.style={
    colback=orange!8,
    colframe=accentorange,
    arc=2mm,
    boxrule=0.8pt,
    left=2mm,
    right=2mm,
    top=1.5mm,
    bottom=1.5mm
  }
}

\title{
  \vspace{-1.5cm}
  \textbf{选题前置调研报告}\\[0.5em]
  \large TODO：项目副标题或英文标题
}

\author{TODO}
\date{\today}

\begin{document}

\maketitle

\begin{tcolorbox}[mybox,title=一句话概括]
TODO：用 2--4 句话概括项目的真正研究对象、目标用户、核心场景、预期贡献和最重要的可行性判断。
\end{tcolorbox}

\vspace{0.5em}

\begin{center}
\textbf{目标投稿方向：TODO}\\
\textbf{关键词：TODO}
\end{center}

\newpage
\tableofcontents
\newpage

\section{项目核心判断}

\subsection{核心问题不是 TODO}

TODO。

\subsection{真正值得研究的问题}

\begin{tcolorbox}[keybox,title=本项目的核心研究问题]
TODO。
\end{tcolorbox}

\subsection{为什么这个问题重要}

\begin{itemize}[leftmargin=2em]
  \item TODO。
  \item TODO。
  \item TODO。
\end{itemize}

\subsection{本项目定位}

\begin{table}[H]
\centering
\caption{本项目定位}
\begin{tabularx}{\textwidth}{p{3cm}X}
\toprule
\textbf{维度} & \textbf{定位} \\
\midrule
不是 & TODO \\
\midrule
是 & TODO \\
\midrule
核心输入 & TODO \\
\midrule
核心输出 & TODO \\
\midrule
核心价值 & TODO \\
\bottomrule
\end{tabularx}
\end{table}

\section{为什么适合目标会议或期刊}

\subsection{目标 venue 的关注点}

TODO。

\subsection{与最接近工作的差异}

\begin{table}[H]
\centering
\caption{与既有工作的差异}
\begin{tabularx}{\textwidth}{p{3.2cm}p{4cm}X}
\toprule
\textbf{维度} & \textbf{既有工作} & \textbf{本项目} \\
\midrule
场景 & TODO & TODO \\
\midrule
用户任务 & TODO & TODO \\
\midrule
系统角色 & TODO & TODO \\
\midrule
输出形式 & TODO & TODO \\
\midrule
研究焦点 & TODO & TODO \\
\bottomrule
\end{tabularx}
\end{table}

\subsection{可能贡献}

\begin{enumerate}[leftmargin=2em]
  \item \textbf{经验贡献：} TODO。
  \item \textbf{概念贡献：} TODO。
  \item \textbf{系统/方法贡献：} TODO。
  \item \textbf{评估贡献：} TODO。
\end{enumerate}

\section{核心概念与研究问题}

\subsection{核心概念定义}

TODO。

\subsection{形式化表达或分析框架}

TODO。

\subsection{研究问题}

\begin{table}[H]
\centering
\caption{研究问题}
\begin{tabularx}{\textwidth}{p{2cm}X}
\toprule
\textbf{编号} & \textbf{研究问题} \\
\midrule
RQ1 & TODO \\
\midrule
RQ2 & TODO \\
\midrule
RQ3 & TODO \\
\midrule
RQ4 & TODO \\
\bottomrule
\end{tabularx}
\end{table}

\section{输入模态、数据模态或研究对象边界}

\subsection{为什么必须先界定边界}

TODO。

\subsection{可选模态或数据来源}

\begin{table}[H]
\centering
\caption{可选模态或数据来源对比}
\begin{tabularx}{\textwidth}{p{3cm}XX}
\toprule
\textbf{方案} & \textbf{优点} & \textbf{局限} \\
\midrule
TODO & TODO & TODO \\
\midrule
TODO & TODO & TODO \\
\midrule
TODO & TODO & TODO \\
\bottomrule
\end{tabularx}
\end{table}

\subsection{本项目建议边界}

\begin{tcolorbox}[warnbox,title=需要对审稿人明确的边界]
TODO。
\end{tcolorbox}

\section{系统方案或方法方案}

\subsection{系统/方法目标}

TODO。

\subsection{输入与输出}

\begin{table}[H]
\centering
\caption{输入与输出}
\begin{tabularx}{\textwidth}{p{3cm}X}
\toprule
\textbf{类别} & \textbf{内容} \\
\midrule
输入 & TODO \\
\midrule
输出 & TODO \\
\bottomrule
\end{tabularx}
\end{table}

\subsection{工作流}

\begin{enumerate}[leftmargin=2em]
  \item TODO。
  \item TODO。
  \item TODO。
  \item TODO。
\end{enumerate}

\section{交互/方法设计原则}

\begin{table}[H]
\centering
\caption{设计原则}
\begin{tabularx}{\textwidth}{p{3cm}p{4cm}X}
\toprule
\textbf{原则} & \textbf{含义} & \textbf{设计落实} \\
\midrule
低打断 / 低负担 & TODO & TODO \\
\midrule
用户主导 & TODO & TODO \\
\midrule
可解释或可追溯 & TODO & TODO \\
\midrule
低置信度退让 & TODO & TODO \\
\bottomrule
\end{tabularx}
\end{table}

\section{算法与数据的最小可行方案}

\subsection{技术路线}

TODO。

\subsection{数据集候选}

\begin{longtable}{p{3cm}p{4cm}p{4cm}p{3cm}}
\caption{数据集候选与可行性} \\
\toprule
\textbf{数据集} & \textbf{可支持内容} & \textbf{不能支持内容} & \textbf{风险} \\
\midrule
\endfirsthead
\toprule
\textbf{数据集} & \textbf{可支持内容} & \textbf{不能支持内容} & \textbf{风险} \\
\midrule
\endhead
TODO & TODO & TODO & TODO \\
\bottomrule
\end{longtable}

\subsection{质量检查}

\begin{itemize}[leftmargin=2em]
  \item TODO。
  \item TODO。
  \item TODO。
\end{itemize}

\section{典型案例或使用场景}

\subsection{案例一：TODO}

\begin{tcolorbox}[mybox,title=Demo Case 1]
\textbf{场景输入：} TODO\\[0.3em]
\textbf{系统/方法判断：} TODO\\[0.3em]
\textbf{推荐输出：} TODO\\[0.3em]
\textbf{评价信号：} TODO\\[0.3em]
\textbf{局限：} TODO
\end{tcolorbox}

\subsection{案例二：TODO}

TODO。

\section{研究问题与评价方案}

\subsection{评价整体设计}

TODO。

\begin{table}[H]
\centering
\caption{评价路线}
\begin{tabularx}{\textwidth}{p{2.8cm}p{3.5cm}p{3.8cm}X}
\toprule
\textbf{阶段} & \textbf{目标} & \textbf{方法} & \textbf{产出} \\
\midrule
Study 1 & TODO & TODO & TODO \\
\midrule
Study 2 & TODO & TODO & TODO \\
\midrule
Study 3 & TODO & TODO & TODO \\
\bottomrule
\end{tabularx}
\end{table}

\subsection{指标}

\begin{table}[H]
\centering
\caption{评价指标}
\begin{tabularx}{\textwidth}{p{4cm}X}
\toprule
\textbf{指标类别} & \textbf{具体指标} \\
\midrule
质量 & TODO \\
\midrule
效率 & TODO \\
\midrule
负担 & TODO \\
\midrule
信任/采纳 & TODO \\
\midrule
结果 & TODO \\
\bottomrule
\end{tabularx}
\end{table}

\section{数据集支持与可行性判断}

\begin{tcolorbox}[keybox,title=可行性结论]
\textbf{Verdict：} TODO：Pass / Weak Pass / Fail。
\end{tcolorbox}

TODO。

\section{标题润色方案与摘要}

\subsection{推荐标题}

\begin{enumerate}[leftmargin=2em]
  \item TODO。
  \item TODO。
  \item TODO。
\end{enumerate}

\subsection{推荐摘要}

TODO。

\section{预期论文结构}

\begin{enumerate}[leftmargin=2em]
  \item Introduction
  \item Related Work
  \item Formative Study / Problem Analysis
  \item Conceptual Framework
  \item System / Method Design
  \item Evaluation
  \item Findings
  \item Discussion
  \item Limitations and Future Work
  \item Conclusion
\end{enumerate}

\section{投稿定位}

\begin{table}[H]
\centering
\caption{投稿目标对比}
\begin{tabularx}{\textwidth}{p{3cm}XX}
\toprule
\textbf{目标} & \textbf{优势} & \textbf{风险} \\
\midrule
TODO & TODO & TODO \\
\midrule
TODO & TODO & TODO \\
\bottomrule
\end{tabularx}
\end{table}

\section{风险与应对}

\begin{longtable}{p{4cm}p{5cm}p{5cm}}
\caption{主要风险与应对策略} \\
\toprule
\textbf{风险} & \textbf{表现} & \textbf{应对} \\
\midrule
\endfirsthead
\toprule
\textbf{风险} & \textbf{表现} & \textbf{应对} \\
\midrule
\endhead
选题太泛 & TODO & TODO \\
\midrule
数据不支撑 & TODO & TODO \\
\midrule
贡献不清 & TODO & TODO \\
\midrule
评估不真实 & TODO & TODO \\
\midrule
伦理或隐私风险 & TODO & TODO \\
\bottomrule
\end{longtable}

\section{下一步工作}

\subsection{近期两周}

\begin{itemize}[leftmargin=2em]
  \item TODO。
\end{itemize}

\subsection{一个月内}

\begin{itemize}[leftmargin=2em]
  \item TODO。
\end{itemize}

\subsection{两到四个月内}

\begin{itemize}[leftmargin=2em]
  \item TODO。
\end{itemize}

\section{最终总结}

\begin{tcolorbox}[keybox,title=最终总结]
TODO：用一段强判断收束，说明这个方向应该继续、收缩还是放弃；下一步最该做什么；哪些话现在不能写进论文。
\end{tcolorbox}

\end{document}
